\subsection{表示的张量积}

我们在第 1 号中定义了 g 的一族表示的直和。现在定义表示的其他运算。

设 \(\mathfrak{g}_1, \mathfrak{g}_2\) 是 \(\mathbf{K}\) 上的两个李代数，\(\mathbf{M}_i\) 是一个 \(\mathfrak{g}_i\)-模（\(i = 1, 2\)）。令 \(\mathbf{U}_i\) 为 \(\mathfrak{g}_i\) 的包络代数，\(\sigma_i\) 为从 \(\mathfrak{g}_i\) 到 \(\mathbf{U}_i\) 的典范映射。于是 \(\mathbf{M}_i\) 是左 \(\mathbf{U}_i\)-模，因而 \(\mathbf{M}_1 \otimes_{\mathbf{K}} \mathbf{M}_2\) 具有典范的左 \((\mathbf{U}_1 \otimes_{\mathbf{K}} \mathbf{U}_2)\)-模结构。现在 \(\mathbf{U}_1 \otimes_{\mathbf{K}} \mathbf{U}_2\) 是 \(\mathfrak{g}_1 \times \mathfrak{g}_2\) 的包络代数，而映射 \((x_1, x_2) \mapsto \sigma_1(x_1) \otimes 1 + 1 \otimes \sigma_2(x_2)\) 是从 \(\mathfrak{g}_1 \times \mathfrak{g}_2\) 到此包络代数的典范映射（§ 2，第 2 号）。因此，在 \((\mathfrak{g}_1 \times \mathfrak{g}_2)\) 上存在 \(\mathbf{M} = \mathbf{M}_1 \otimes_{\mathbf{K}} \mathbf{M}_2\)-模结构，使得：

\[
\begin{array}{r l} (x _ {1}, x _ {2}) _ {\mathrm{M}}. (m _ {1} \otimes m _ {2}) & = (\sigma_ {1} (x _ {1}) \otimes 1 + 1 \otimes \sigma_ {2} (x _ {2})) . (m _ {1} \otimes m _ {2}) \\ & = ((x _ {1}) _ {\mathrm{M} _ {1}}. m _ {1}) \otimes m _ {2} + m _ {1} \otimes ((x _ {2}) _ {\mathrm{M} _ {2}}. m _ {2}). \end{array}\tag{1}
\]

此结构定义了 \(\mathfrak{g}_1\times \mathfrak{g}_2\) 在 \(\mathbf{M}\) 上的表示。

若现在 \(\mathfrak{g}_1 = \mathfrak{g}_2 = \mathfrak{g}\)，则从 \(x \mapsto (x, x)\) 到 \(\mathfrak{g}\) 的同态 \(\mathfrak{g} \times \mathfrak{g}\) 与上述表示复合，定义了 \(\mathfrak{g}\) 在 \(\mathbf{M}\) 上的表示，因而在 \(\mathfrak{g}\) 上定义了 \(\mathbf{M}\)-模结构，使得：

\[
x _ {\mathrm{M}}. (m _ {1} \otimes m _ {2}) = (x _ {\mathrm{M} _ {1}}. m _ {1}) \otimes m _ {2} + m _ {1} \otimes (x _ {\mathrm{M} _ {2}}. m _ {2}).\tag{2}
\]

同理可见：

命题 2. 设 g 是 K 上的李代数，\(M_{i}\) 是 g-模（\(1 \leqslant i \leqslant n\)）。在张量积 \(M_{1} \otimes_{K} M_{2} \otimes \cdots \otimes M_{n}\) 上存在唯一的 g-模结构，使得

\[
x _ {\mathrm{M}}. (m _ {1} \otimes \dots \otimes m _ {n}) = \sum_ {i = 1} ^ {n} m _ {1} \otimes \dots \otimes (x _ {\mathrm{M} _ {i}}. m _ {i}) \otimes \dots \otimes m _ {n}\tag{3}
\]

对所有 \(x \in \mathfrak{g}\)、\(m_1 \in \mathbf{M}_1, \ldots, m_n \in \mathbf{M}_n\)。

相应的表示称为给定的 g 在各 \(M_{i}\) 上的表示的张量积。

特别地，若 \(\mathbf{M}\) 是 \(\mathfrak{g}\)-模，命题 2 定义了 \(\mathfrak{g}\)-模结构
在每个 \(\mathbf{M}_p = \bigotimes \mathbf{M}\) 上定义了结构，因而也在 \(\mathbf{T}\) 的张量代数 \(\mathbf{M}\) 上定义了结构。

公式 (3) 表明，对所有 \(x \in \mathfrak{g}\)，\(x_{\mathrm{T}}\) 是代数 \(\mathrm{T}\) 上延拓 \(x_{\mathrm{M}}\) 的唯一导子。我们知道（§ 2，第 8 号），\(x_{\mathrm{T}}\) 在取商后在 \(\mathrm{S}\) 的对称代数 \(\mathrm{M}\) 上定义了一个导子。因此，可以把 \(\mathrm{S}\) 看作 \(\mathfrak{g}\) 的商 \(\mathrm{T}\)-模，而 \(x_{\mathrm{S}}\) 是 \(\mathrm{S}\) 的导子。

更具体地，把 g 借助其伴随表示看作 g-模。设 U 是 g 的包络代数。根据 §2 的命题 7，\(x_{\mathrm{M}}\) 在取商后定义了 U 上的导子，而这正是由 \(\sigma(x)\) 定义的内导子（\(\sigma\) 表示从 g 到 U 的典范映射）。于是可以把 U 看作 T 的商 g-模。如果 K 是特征为 0 的域，则从 S 到 U 的典范同构是 g-模同构（§2，第 8 号）。

